\section{Gap curvature from one upper time}\label{sec:anchored-curvature}

The exterior at \(T\) gives the exact initial slope in \eqref{eq:exterior-gap}.
To control the error of extrapolating this slope, tilt the complete law at
\(T\) by \(e^{-sU^2}\), as in \eqref{eq:whole}. Throughout, the saddle
parameters \(L,\kappa,S\) of \eqref{eq:saddle} are evaluated at \(T\).
Set
\[
 \lambda=L\sqrt\kappa,\quad r=\tfrac12\sqrt{\log\kappa},\quad
 \gamma=2L/\sqrt\kappa,\quad b=\gamma s,\quad
 D_*=\pi^2/(4L^2),\quad \eta=\kappa^{-3/4} .
\]
Here \(s\) is elapsed backward physical time, and \(y=Lx\) is the scaled
spatial argument. A joint complex estimate in \(b\) and \(y\) controls the
difference of the two inverse positions before time differentiation
\cite{EZCurvature}.

\Needspace{8\baselineskip}
\begin{theorem}[Signed remainder on a backward interval]\label{thm:anchored-curvature}
For all sufficiently large \(T\), the marked zeros for each \(k\ge1\) with
\(q_{k+1}\le\lambda r/8\) stay simple, adjacent and separated from the rest
of the divisor throughout \(0\le s\le s_*:=\gamma^{-1}\). Uniformly over
these pairs and times, with primes denoting derivatives in physical time,
\begin{align}
 \frac{D_k'(T-s)}{D_*}
   &=-\frac4\kappa\{1+b/\lambda+6\alpha b+O(\eta)\}<0,
      \qquad \alpha=\frac{2S}{3\kappa^{3/2}},\label{eq:anchored-slope}\\
 \frac{D_k''(T-s)}{D_*}
   &=\frac8{\kappa^2}+\frac{32LS}{\kappa^3}+O(L\kappa^{-9/4})
     \sim\frac{32L}{\kappa^2}.\label{eq:anchored-curvature}
\end{align}
In particular, the exact anchor slope \(E_k=D_k'(T)=2-g_k(T)^2M_k(T)\)
gives the signed estimate
\begin{equation}\label{eq:anchored-linear}
 \frac{D_k(T-s)-D_k(T)+sE_k}{D_*}
   =s^2\left\{\frac4{\kappa^2}+\frac{16LS}{\kappa^3}
                  +O(L\kappa^{-9/4})\right\} .
\end{equation}
It is strictly positive for every \(0<s\le s_*\).
For each Taylor degree \(m\ge2\) the bound with exact coefficients is
\begin{equation}\label{eq:anchored-taylor}
 \frac{\bigl|D_k(T-s)-\sum_{j=0}^m(-s)^j D_k^{(j)}(T)/j!\bigr|}{D_*}
 \le \frac{C\eta}{\lambda}\,\frac{(|b|/2)^{m+1}}{1-|b|/2}
 \qquad (0\le s\le s_*).
\end{equation}
The constant \(C\) is independent of the pair, Taylor degree and partial time.
\end{theorem}

\begin{proof}
\emph{The joint transform.}
Write \(W=\sqrt\kappa(U-L)\) and \(\zeta=y/\lambda\). Removing the scalar
\(e^{-sL^2}\) leaves the tilt \(e^{-bW-bW^2/(2\lambda)}\).
Use the unnormalized density \(F\) of Lemma~\ref{lem:saddle}, extended by
zero for \(w<-\lambda\), and the Gaussian model
\(\varphi(w)=e^{-w^2/2}/\sqrt{2\pi}\). Define
\[
 a_\pm=-b\pm\ii\zeta,\qquad v_b=A_b^{-1},\qquad A_b=1+b/\lambda,
\]
\[
 J_\pm=\int_{\mathbb R} F(w)e^{a_\pm w-bw^2/(2\lambda)}\dd w,
 \qquad I_\pm=v_b^{1/2}e^{v_ba_\pm^2/2},\qquad m_\pm=v_b a_\pm .
\]
Up to the removed scalar and a common nonzero anchor constant, the field is
\((e^{\ii y}J_++e^{-\ii y}J_-)/2\).
The integral with \(\varphi\) in place of \(F\) is \(I_\pm\), with its square
root continued from \(v_b=1\). Three derivatives in \(a_\pm\) show that
the Gaussian integral with the cubic correction in Lemma~\ref{lem:saddle}
is \(I_\pm\{1-\alpha(m_\pm^3+3v_bm_\pm)\}\).

For \(|b|\le3,\ |\zeta|\le r\), complete the real square in the absolute
central error integral on \(|w|\le w_0=\kappa^{1/12}\). Its tilted polynomial
moment is \(O(1+(r+3)^6)\). Division by \(|I_\pm|\) leaves a bounded
prefactor and the exponential of
\[
 \frac{(\Real a_\pm)^2}{2\Real A_b}-\frac12\Real(v_ba_\pm^2)
 =\frac12(\Imag a_\pm)^2+O((r+3)^2/\lambda).
\]
The resulting relative central error is \(O(\eta)\), since the density
error supplies \(\kappa^{-1}\) and the exponential costs at most
\(\kappa^{1/8+o(1)}\). The complete source on \(w\ge-\sqrt\kappa/4\)
outside the central interval is bounded by \(Ce^{-cw^2}\). The tilt
costs at most \(e^{(r+3)|w|+3w^2/(2\lambda)}\), leaving a fixed Gaussian
exponent beyond \(w_0\). The Gaussian model tails obey the same estimate,
including beyond the source's left endpoint. On
\(-\lambda\le w<-\sqrt\kappa/4\) the source is \(O(e^{-c\kappa})\) and
the integrated tilt costs at most
\(\lambda e^{(r+3)\lambda+3\lambda/2}=e^{o(\kappa)}\).
Thus all tail errors are smaller than any inverse power of \(\kappa\)
even after division by \(|I_\pm|\).

It follows that \(J_\pm/I_\pm=1-\alpha(m_\pm^3+3v_bm_\pm)+O(\eta)\). The deviation
from one is \(o(1)\) and its square is \(o(\eta)\), so
\begin{equation}\label{eq:anchored-log}
 \log J_\pm=\tfrac12\log v_b+\tfrac12 v_ba_\pm^2
                  -\alpha(m_\pm^3+3v_bm_\pm)+O(\eta).
\end{equation}
The integrals are holomorphic jointly by the same dominating bounds on
the complex disks. The remainder in \eqref{eq:anchored-log} is therefore
jointly holomorphic in the interiors and uniformly bounded by \(C\eta\).

\emph{The counted phase and its inverse.}
The field factors as a nonvanishing amplitude times \(\cos\Theta_b(y)\), where
\begin{equation}\label{eq:anchored-phase}
 \Theta_b(y)=y+\frac{\log J_+-\log J_-}{2\ii}
 =v_b y-\alpha\{(3v_b^3 b^2+3v_b^2)\zeta-v_b^3\zeta^3\}+O(\eta).
\end{equation}
The Gaussian part uses \(1-v_b b/\lambda=v_b\). The phase is odd in
\(y\), by exchange of \(J_+\) and \(J_-\), and real for real \(b,y\).
In the coordinate \(p=v_by\), the correction to \(p\) has size
\(O(\alpha(1+r^3)+\eta)\) on \(|p|\le3\lambda r/4\).
Rouché on \(|p|=\lambda r/2\) gives a unique simple inverse
\(y=Y_b(q)\) for \(|q|\le\lambda r/4\), jointly holomorphic with room at
the \(q\)-boundary. Cauchy's estimate on the larger \(p\)-disk gives
slope \(1+o(1)\) in \(p\) on the real segment through these inverses.
Thus \(\Theta_b(0)=0\) and the factorization count every positive zero
through \(Y_b(\lambda r/4)/L\) for real \(0\le b\le1\). In particular
the marked endpoints are \(Y_b(q_k)/L,Y_b(q_{k+1})/L\).

The inverse equation first gives \(p-q=O(\alpha(1+r^3)+\eta)\).
Replacing \(p\) by \(q\) in the cubic correction costs at most
\(O(\alpha(1+r^2)|p-q|/\lambda)=o(\eta)\). Replacing \(v_b\) by one
in that correction costs \(O(\alpha(1+r^3)/\lambda)=o(\eta)\). Hence
\begin{equation}\label{eq:anchored-inverse}
 Y_b(q)=q/v_b+\alpha\{(3b^2+3)q/\lambda-(q/\lambda)^3\}+R_Y(b,q),
 \qquad |R_Y|\le C\eta .
\end{equation}
Cauchy in \(q\) gives \(|\partial_q R_Y|\le C\eta/(\lambda r)\) on
\([0,\lambda r/8]\), uniformly in the time disk. Subtract
\eqref{eq:anchored-inverse} at the two adjacent levels before squaring.
In units of \(D_*\) the resulting holomorphic squared gap is
\begin{align}
 \mathcal D_k(b)&=\left(\frac{Y_b(q_{k+1})-Y_b(q_k)}{\pi}\right)^2
                   =P_k(b)+R_k(b),\label{eq:anchored-model}\\
 P_k(b)&=(1+b/\lambda)^2+\frac{2\alpha}{\lambda}
      \left(3b^2+3-\frac{q_{k+1}^3-q_k^3}{\pi\lambda^2}\right),
      \qquad |R_k(b)|\le C\eta/\lambda\quad (|b|\le3).\notag
\end{align}
The cross correction from replacing \(1/v_b\) by one in the order-\(\alpha\)
term after squaring costs \(O(\alpha(1+r^2)/\lambda^2)\). The square
of the gap correction costs \(O(\alpha^2(1+r^2)^2/\lambda^2)\).
Both fit the remainder along with the adjacent difference of \(R_Y\).

\emph{Physical time and clearance.}
For real \(s\) the squared gap satisfies \(D_k(T-s)=D_*\mathcal D_k(b)\).
The first two \(b\)-derivatives of \(R_k\) are \(O(\eta/\lambda)\) on
\(0\le b\le1\) by Cauchy. Using \(\mathrm d/\mathrm dt=-\gamma\,\mathrm d/\mathrm db\)
gives \eqref{eq:anchored-slope} and \eqref{eq:anchored-curvature}.
The two curvature terms come from the Gaussian part and the cubic source
correction, respectively. Their sum also equals the relative curvature
at \(T\) of the instantaneous lattice value \(\pi^2/(4L_t^2)\), by
\(L_t'=2L_t/\kappa_t\) and \eqref{eq:saddle}. Integration against \((s-u)\dd u\) for
\(0\le u\le s\) proves \eqref{eq:anchored-linear}. Subtracting the degree-\(m\)
Taylor polynomial of \(\mathcal D_k\) leaves only the Taylor tail of \(R_k\)
when \(m\ge2\). Its radius-two Cauchy bound gives \eqref{eq:anchored-taylor}.

The slope and \eqref{eq:anchored-model} give \(D_k(T-s)\ge D_k(T)>D_*/2\).
The outer level \(q_{k+2}\) lies within \(\lambda r/4\) eventually.
The next zero on the left is counted there as well, with the reflected
first zero used for \(k=1\). The uniform \(o(1)\) displacement from \(q/v_b\)
in \eqref{eq:anchored-inverse} gives adjacent spacings
\((\pi/L)(1+o(1))\). The field is all real at these large physical times
\cite{Polymath2019}, so a circle of radius \(g_k\) about the real midpoint
contains exactly the pair and stays clear of the other zeros.
\end{proof}

\subsection{The anchor coefficients and their precision}
The complete hierarchy identifies the derivatives in
\eqref{eq:anchored-taylor} on the separated chart. With the pair index
suppressed, combine \eqref{eq:exterior-gap} with the diagonal equation
in \eqref{eq:parent} to give
\begin{equation}\label{eq:anchored-polynomial}
 D''=-8M+8DM^2+48D S_{2,2}+64D^2 S_{3,3} .
\end{equation}
The unreduced integer split recurrence raises the largest total pole
order of a row by at most two per differentiation. Applied to \(D'=2-4DM\),
it uses rows of order at most \(2m\) for the Taylor jet of degree \(m\).
The diagonal reduction in \eqref{eq:anchored-polynomial} uses an order-six
row \cite[\S3]{AGC}. The joint disk bounds the resulting scalar gap. An
enclosure for an individual evolving row requires its own analytic estimate.

For the finite anchor calculation use \(y_i=Lx_i(T)\) and take the head
through \(q_K\le\lambda r/4\) with \(K\) maximal. The marked endpoints have
an upper bound \(A\) and the unknown positive roots a lower bound \(R\),
with \(R\asymp K\) and \(A/R\le2/3\) eventually by \eqref{eq:anchored-inverse}.
Apply \eqref{eq:exterior-tail} in these coordinates, with
\(W_j^{(y)}=\sum_{i>K} y_i^{-2j}\) and moments \(\nu_j=\E_T(U/L)^j\).
For any fixed row of order \(h=p+q\), retaining degrees \(d<N\) uses
moments through order at most \(h+N-1\). Each known exterior occurrence
in that row is evaluated with positive lower distances \(l_+,l_-\).
For position errors \(e_i,e_+,e_-\) its error is at most
\[
 \frac{p(e_i+e_+)}{l_+^{p+1}l_-^q}
   +\frac{q(e_i+e_-)}{l_+^p l_-^{q+1}},
\]
with terms of zero numerator omitted. The same endpoint enclosures enter
each row and the gap. Widths from the radial-power subtractions and
the target coefficients in the retained tail must be propagated as well.

Here is a sufficient acquisition budget for the slope and curvature
\cite[finite-precision supplement]{EZCurvature}. In the scaled variables
\(D_y=L^2D\), \(M_y=M/L^2\) and \(S_{p,q}^{(y)}=S_{p,q}/L^{p+q}\), use
\(\sigma=5/6\) in the bound following \eqref{eq:exterior-tail} and choose
\(N=O(\log\lambda)\) with \((4/5)^N\le\lambda^{-6}\). The omitted tail
at orders \(h=2,4,6\) is \(O(\lambda^{-6})\), since
\(W_1^{(y)}\le\nu_2/2=O(1)\) and \(R^{2-h}\le1\) eventually. The retained
degrees need moments through order \(N+5\).

On the unit disk the truncated radial cosine polynomial lies within
\(3/4\) of one, with strict clearance for sufficiently accurate input
enclosures, since the corresponding full positive coefficient sum tends
to \(\cosh 1-1\). This limit follows by the weighted saddle bounds. If
\(\delta\) is the sum of absolute errors in its normalized coefficients,
Cauchy's estimate for the perturbed logarithm gives an error
\(O(j\delta)\) in the full \(j\)-th radial power. A common position radius
\(e_{\rm pos}\) for the scaled head gives an additional \(O(je_{\rm pos})\)
in the head subtracted from it, using \(y_i-e_{\rm pos}\ge i\) and summing
the derivative bounds for \(y_i^{-2j}\). The retained coefficients satisfy
\(|h_d|\le\binom{h+d-1}{d}A^d\), so the summed error from these radial
subtractions is at most
\begin{equation}\label{eq:anchored-input-width}
 O\big((\delta+e_{\rm pos})N^7\max(1,A)^N\big).
\end{equation}
With the exact radial powers held fixed, the endpoint errors in the tail
cost \(O(e_{\rm pos})\) by differentiating the convergent reflected sums
under the radial clearance. The counted exterior head has the same bound
from its separated spacings and summable reciprocal products.

Thus \(\delta,e_{\rm pos}\ll\lambda^{-6}/(N^7\max(1,A)^N)\) suffice,
which calls for \(O(\log^2\lambda)\) absolute bits in these scaled data.
The inputs of the polynomial for \(D''/L^2\) in
\eqref{eq:anchored-polynomial} are bounded in the scaled coordinates.
Errors \(O(\lambda^{-6})\) persist at that order in the polynomial and in
\(E_k=2-(Lg_k)^2M_y\). The curvature signal at the anchor is
\(D''/L^2\sim8\pi^2L/\lambda^4\), which dominates this error.

The full odd phase also gives the head to the required precision. In
\(y\), its linear coefficient is \(a/L\). The remaining odd part is bounded
by \(M_{\rm ph}=O(\alpha(1+r^3)+\eta)\) on \(|y|\le\lambda r\), by
\eqref{eq:anchored-phase} at \(b=0\) and Cauchy for the linear term.
On \(|y|\le X:=\lambda r/2\), its Taylor tail past degree \(P\) is
\(O(M_{\rm ph}2^{-P})\). Take \(P=O(\log^2\lambda)\) to make the tail
smaller than \(e_{\rm pos}\). The coefficients are obtained from the log
of the truncated moment generating polynomial for \(U/L\). On the disk
of radius \(1/4\), the positive coefficient sum excluding the constant
tends to \(e^{1/4}-1<1/2\) when all moments are included. Thus an error
\(\delta_{\rm ph}\) in the sum of coefficient radii weighted at \(1/4\)
costs \(O(4^j\delta_{\rm ph})\) in the \(j\)-th log coefficient. It suffices
that \(\delta_{\rm ph}\ll e_{\rm pos}(4X)^{-P}\). The true phase slope
is bounded below and the counted head lies strictly inside \([0,X]\), so
the resulting \(O(e_{\rm pos})\) phase error bounds the positional error.
This gives a sufficient moment order \(O(\log^2\lambda)\) and
\(O(\log^3\lambda)\) absolute bits in the stated weighted coefficient norm,
including acquisition of the head. These are asymptotic tolerances.
An application at a specified \(T\) requires checking the inequalities
with the actual input enclosures.

For a finite Taylor calculation write \(d_j\) for the absolute error
in \((-1)^j D_k^{(j)}(T)/(D_*\gamma^j j!)\), allowing for scale uncertainty
and rounding. Add \(\sum_{j=0}^m d_j|b|^j\) to \eqref{eq:anchored-taylor}.
For example, \(\sum_{j=2}^m d_j=O(\eta/\lambda)\) adds a normalized
error \(O((\eta/\lambda)b^2)\) for \(|b|\le1\).
In \eqref{eq:anchored-linear} inexact anchor values add physical errors
\(|\delta D_k(T)|+s|\delta E_k|\). Resolving its signed quadratic increment
at a specified \(b>0\) requires \(d_0+d_1|b|\) small compared with
\((\alpha/\lambda)b^2\). Exact shared anchor values yield the \(s^2\)
vanishing down to \(s=0\). The finite coefficient budget is for the
anchor data, while \eqref{eq:anchored-model} supplies the complete analytic
time allowance from the source.

\subsection{Count, scope and comparison}
The proof counts \(K=\lfloor\lambda r/(4\pi)+1/2\rfloor\) positive zeros
through \(Y_b(\lambda r/4)/L\). The retained pairs number
\(\lfloor\lambda r/(8\pi)+1/2\rfloor-1\), of order
\(\sqrt{T}(\log T)^2\), on a physical interval of length
\(s_*\sim2\sqrt{2T/\log T}=o(T)\). The onset in \(T\) is unevaluated.
The current estimate of Section~\ref{sec:transport} treats a fixed multiple
of the saddle window using bounds along the interval. The present estimate
uses one upper law and a time disk on a growing window. A separate spatial
estimate treats sufficiently high labels on a fixed slab in \((0,1/2]\),
with its own direct two-time comparison \cite[\S5 and references there]{EZCurvature}.
A backward path from the present upper anchor to a fixed lower time
requires different bounds.

Direct evolution by \eqref{eq:whole} uses the same complete law. Integrating
\(\Theta_b'\) between the two counted roots gives \(\pi\), so the direct
phase calculation with the same spatial Cauchy estimate and time disk
recovers \eqref{eq:anchored-model} and the same curvature and Taylor
bounds. The shared analytic error provides the gain over separate endpoint
intervals, whose widths can persist as \(s\) decreases to zero.
The anchor estimate identifies these bounds with the derivatives from
the exterior hierarchy.

The sign uses the quantitative tilted theta expansion. For example, a
stationary cosine has zero gap curvature and its atomic law lacks the cubic
saddle correction. Arithmetic and profile modifications change the complete
source, and changed-observation controls fail the literal heat-cosine tilt.
The tail identity \eqref{eq:exterior-tail} requires a complete all-real
divisor with integer multiplicities. Controls with nonreal zeros or relaxed
divisor weights fail this premise. Finally, a backward collision of the
atomic field \(e^t\cos x+c\) destroys the separated chart
\cite[\S5 and references there]{EZCurvature}.
